\documentclass[11pt]{article}

\usepackage[english]{babel}
\usepackage[utf8]{inputenc}
\usepackage{amsmath}
\usepackage{amssymb}
\usepackage{graphicx}
\usepackage[colorinlistoftodos]{todonotes}
\usepackage{listings,multicol}
\usepackage{textcomp}
\usepackage{hyperref}

\setlength{\oddsidemargin}{0.5cm} \setlength{\evensidemargin}{0cm}
\setlength{\textwidth}{16cm} \setlength{\textheight}{23cm}
\setlength{\topmargin}{-0.5cm}
\textheight 21.5cm


\begin{document}

\title{G3 NUMERICO 220023-220138 UBB}

{\begin{minipage}{2cm}
\hspace*{1cm}\includegraphics[width=0.6\textwidth]{escubo-ubb.eps}
\end{minipage}
\begin{minipage}{12cm}
\small
{\bf \rm 
{
\begin{center}
{\footnotesize UNIVERSIDAD DEL B\'IO-B\'IO} \\
{\scriptsize FACULTAD DE CIENCIAS}  \\
{\scriptsize DEPARTAMENTO DE MATEM\'ATICA}  \\
{\scriptsize Profesor:  Franco Milanese}\\
{\scriptsize Primer Semestre de 2016}
\end{center}
}}
\end{minipage}}
{\begin{minipage}{2cm}
\hspace*{-0.5cm}\vspace*{-0.05cm}\includegraphics[width=0.7\textwidth]{escudo-dmat.eps}
\end{minipage}}

\hspace*{-1,5cm}\rotatebox[origin=c]{90}{\begin{picture}(0,0)
\put(0,7){\makebox(9,-13)[l]{\hspace*{-6.5in} \bf \it Departamento de Matem\'atica - Universidad del B\'io-B\'io - 2016}}
\end{picture}}

\vspace*{0.5cm} \centerline {\bf\underline{Gu\'ia 3, M\'etodos Num\'ericos I 220023-220138 }}
\centerline{\textrm{.}}  \vspace{0.2cm}


% \textbf{Nombre:} \hspace{0.5\textwidth}\textbf{Carrera:}
% \vspace{0.1cm}
% \textbf{Profesor:}\hspace{0.5\textwidth} \textbf{ RUT:}
%  \begin{center}
%  \begin{tabular}{||p{2cm}|p{2cm}|p{2cm}||}
%  \hline
%  Pregunta 1 &  Pregunta 2 &     Total\\
%  \hline

%   \vspace{1.5cm} & &       \\
%  \hline
%  \end{tabular}
%  \end{center}
% Enviar documentos solicitados en el formato solicitado a \textbf{veranonumerico@gmail.com}.



\begin{enumerate}

\item Dados los tres puntos $\{(-1,1),(0,0),(1,1)\}$ determine el polinomio interpolante de grado dos,
\begin{itemize}
 \item[a)] usando la base can\'onica,
 \item[b)] usando la base de polinomios de Lagrange,
% \item[c)] usando la base de Newton (diferencias divididas).
\end{itemize}

\item Cuantas multiplicaciones se requieren para evaluar un polinomio $p(t)$ 
de grado $n-1$ en un punto cualquiera $t$ 
\begin{itemize}
 \item[a)] en su representaci\'on respecto a la base can\'onica,
 \item[b)] en su representaci\'on respecto a la base de polinomios de Lagrange,
 %\item[c)] en su representaci\'on de diferencias definidas.
\end{itemize}

\item Ocupa la cota de error para los polinomios interpolantes, para estimar una 
el error m\'aximo al interpolar la funci\'on $sin(x)$ por polinomios de grado $4$ usando 
$5$ puntos equiespaciados en el intervalo $[0,\pi]$. \textquestiondown Cu\'antos puntos son 
necesarios para tener una precisi\'on de $10^{-10}$?.

\item Para un conjunto de puntos de interpolaci\'on $\{t_1,\cdots,t_n\}$, se define la funci\'on
$$
\pi(t)=(t-t_1)(t-t_2)\cdots (t-t_n),
$$
\begin{itemize}
 \item[a)] Pruebe que $\pi'(t_j)=(t_j-t_1)\cdots (t_j-t_{j-1})(t_j-t_{j+1})\cdots (t_j-t_n)$,
 \item[b)] Use lo probado en a) para mostrar que
 $$
 \ell_j(t)=\frac{\pi(t)}{(t-t_j)\pi'(t_j)},
 $$
 donde $\ell_j$ es el $j-$\'esimo polinomio de Lagrange de los puntos $\{t_i\}_{i=1}^n$.
\end{itemize}

\item
Calcule por m\'inimos cuadrados un interpolante de grado $2$ de los puntos $\{(0,1),(1,1),(2,0)\}$.

\item
Calcule por m\'inimos cuadrados la recta que mejor se ajusta a los puntos
$$
\begin{array}{c||c|c|c|c|c|c|c|}
 x	& 8 & 2 & 11 & 6 & 5 & 4
 \\
 \hline
 y	& 3 & 10&  3 & 6 & 8 & 12
 \end{array}$$

\end{enumerate}


\end{document}  